\subsection{Summary of the work completed}
A handwriting robot was designed and built to recognise photographed handwriting, identify a known writer and reproduce new text in that writer's style. Page conditioning and line segmentation feed the text and author classifiers. Labelled samples supply stored glyph trajectories, which are arranged into G-code and executed by the calibrated Cartesian gantry. Eleven profiles were demonstrated in physical output.

\subsection{Summary of observations and findings}
Original-line validation gives 94.06\% character accuracy and 98.08\% author identification. Physical output gives 84.41\% character recovery and 86.41\% line identification, with all fourteen page pools identified correctly. The separate geometric comparison retains 86.68\% of measured natural agreement on 83 repeated words. This supports useful style retention while showing that recognisable identity does not imply identical letter geometry.

All 25 completed ODROID drawing requests meet four seconds per character. All 70 measured feature lengths are within 0.5 mm, with a largest error of 0.49 mm. Complete classification remains outside its time allowance on the recorded ODROID and laptop tests. A measured 2.76-fold network acceleration provides a credible improvement route, but complete embedded compliance requires further measurement and image-processing work.

\subsection{Contribution}
The project demonstrates a complete path from static handwriting samples to physical output in multiple styles. Its first-principles network implementation, explicit glyph geometry and direct motor control make the processing steps inspectable. Classifier scores, natural-variation comparisons and physical dimensions provide complementary evidence of the resulting performance.

\subsection{Suggestions for future work}
The next priorities are output-preserving batching and compiled processing, followed by complete ODROID timing tests. Contextual glyphs and multi-character units should improve connected styles. Repeated bed coordinates and word overlays would extend dimensional verification to placement and repeatability. Improved motor protection, acceleration control and pen restraint would strengthen the next hardware revision.
